% Tradução brasileira revista a partir do workpass espanhol congelado;
% a fonte permanece em source_es.
% SGA 3, Exposé VIII, tradução brasileira: §2.
% Autoridade: Polo--Gille Exp8-8nov09.pdf, pp. locais 7--8,
% leitor completo pp. 623--624.

\setcounter{footnote}{10}

\SGARefTarget{sga3:viii:section:2}\section*{2. Propriedades esquemáticas dos grupos diagonalizáveis}
\addcontentsline{toc}{section}{2. Propriedades esquemáticas dos grupos diagonalizáveis}

Essas propriedades estão resumidas na
\SGARefLink{sga3:viii:proposition:2:1}{proposição seguinte}.

\medskip
\noindent\SGARefTarget{sga3:viii:proposition:2:1}\textbf{Proposição 2.1.}
\textit{Seja \(S\) um pré-esquema não vazio, seja \(M\) um grupo comutativo
ordinário e seja \(G=D(M_S)\) o \(S\)-grupo diagonalizável definido por
\(M\). Então:}
\begin{enumerate}
  \item[(a)]\SGARefTarget{sga3:viii:proposition:2:1:item:a:01}
  \textit{\(G\) é fielmente plano sobre \(S\) e afim sobre \(S\) (a
  fortiori, quase compacto sobre \(S\));}

  \item[(b)]\SGARefTarget{sga3:viii:proposition:2:1:item:b:01}
  \textit{\(M\) é de tipo finito se, e somente se, \(G\) é de tipo finito
  sobre \(S\), se, e somente se, \(G\) é de apresentação finita sobre \(S\);}

  \item[(c)]\SGARefTarget{sga3:viii:proposition:2:1:item:c:01}
  \textit{\(M\) é finito se, e somente se, \(G\) é finito sobre \(S\), se, e
  somente se, \(G\) é de tipo finito sobre \(S\) e é anulado por um inteiro
  \(n>0\). Nesse caso,}
  \[
    \deg(G/S)=\operatorname{Card}(M);
  \]

  \item[\((c')\)]\SGARefTarget{sga3:viii:proposition:2:1:item:cprime:01}
  \textit{\(M\) é um grupo de torção se, e somente se, \(G\) é inteiro
  sobre \(S\);}

  \item[(d)]\SGARefTarget{sga3:viii:proposition:2:1:item:d:01}
  \textit{\(M=0\) se, e somente se, \(G\) é o \(S\)-grupo unidade;}

  \item[(e)]\SGARefTarget{sga3:viii:proposition:2:1:item:e:01}
  \textit{\(M\) é de tipo finito e a ordem de seu subgrupo de torção é
  prima às características residuais de \(S\) se, e somente se, \(G\) é
  liso sobre \(S\).}
\end{enumerate}

A verificação de \SGARefLink{sga3:viii:proposition:2:1:item:a:01}{(a)} a
\SGARefLink{sga3:viii:proposition:2:1:item:d:01}{(d)} é imediata e fica a cargo
do leitor. Demonstremos \SGARefLink{sga3:viii:proposition:2:1:item:e:01}{(e)}.
Se \(G\) é liso sobre \(S\), é localmente de apresentação finita sobre
\(S\), e portanto de apresentação finita sobre \(S\), pois é afim
sobre \(S\); assim, \(M\) é de tipo finito. Podemos supor desde o
início que \(M\) seja de tipo finito e, por conseguinte, que \(G\) seja de
apresentação finita sobre \(S\). Então\footnote{Nota do editor: aqui se
usa que \(G\) é plano sobre \(S\), por
\SGARefLink{sga3:viii:proposition:2:1:item:a:01}{(a)}.} \(G\) é liso sobre
\(S\) se, e somente se, suas fibras geométricas são lisas, o que nos reduz ao
caso em que \(S\) é o espectro de um corpo algebricamente fechado
\(k\).

Escrevamos \(M=T\times L\), onde \(T\) é o subgrupo de torção e \(L\)
é livre, com \(L\simeq\mathbb Z^r\). Então
\[
  D(M)=D(T)\times D(L),
  \qquad
  D(L)\simeq\mathbb G_m^r,
\]
e \(D(L)\) é liso sobre \(k\). Portanto, \(G=D(M)\) é liso sobre \(k\)
se, e somente se, \(D(T)\) é liso. Como \(D(T)\) é finito sobre \(k\)
de grau igual à ordem \(n\) de \(T\), isso significa que \(D(T)(k)\)
tem \(n\) elementos. Ora, \(T\) é isomorfo a uma soma de grupos
\(\mathbb Z/n_i\mathbb Z\), onde \(n=\prod_i n_i\); por conseguinte,
\[
  D(T)=\prod_iD(\mathbb Z/n_i\mathbb Z)
      =\prod_i\mu_{n_i},
\]
onde \(\mu_{n_i}\) é o esquema em grupos das raízes \(n_i\)-ésimas
da unidade. Tem-se
\[
  \operatorname{card}\bigl(D(T)(k)\bigr)
  =\prod_i\operatorname{card}\bigl(\sisi{\mu_i(k)}{\mu_{n_i}(k)}\bigr).
\]
Aqui \(\operatorname{card}(\sisi{\mu_i(k)}{\mu_{n_i}(k)})\), o número de raízes
\(n_i\)-ésimas da unidade em \(k\), é no máximo \(n_i\), e há igualdade
se, e somente se, \(n_i\) é primo à característica \(p\) de \(k\).
% SGA-SRC-000078, ACEITO: original \mu_i(k), atual \mu_{n_i}(k).
Segue-se que
\[
  \operatorname{card}\bigl(D(T)(k)\bigr)=n,
  \qquad n=\prod_i n_i,
\]
se, e somente se, cada \(n_i\) é primo a \(p\), isto é, se, e somente se, \(n\)
é primo a \(p\). \hfill\(\square\)
